\section{Value and Limits of \texorpdfstring{\(\Phi\)}{Phi}}

The value of \(\PhiFCIG\) is operational.  In the source-locked audits, it records finite structure not visible from coarser summaries alone.  The Type-A, modular centered, signed Type-B/C, dihedral, length-colored \(G_2\), projective \(F_4\), projective \(H_3\), projective \(H_4\), projective \(D_5\), projective \(E_6\), projective \(E_7\), projective \(E_8\), and centered-operator artifacts all provide finite evidence that this channel is useful for comparison.

\begin{table}[!htbp]
\centering
\scriptsize
\setlength{\tabcolsep}{3pt}
\renewcommand{\arraystretch}{0.86}
\caption{What the worked witnesses show, I.}
\label{tab:phi-payoff}
\begin{tabular}{@{}L{0.24\textwidth}L{0.31\textwidth}L{0.35\textwidth}@{}}
\toprule
source & finite observation & allowed conclusion \\
\midrule
Type-A Fourier witness & same class histogram; ranks 2 and 1 at \((3,1)\) & \(\Phi\) is not determined by class histograms \\
modular centered witness & \(I_7(1,1)\), \(p=2\), \(\SNFL=[4,4,8]\), rank \(3\to0\) & rational centered rank does not determine modular rank \\
signed Type-B/C witness & same signed-cycle histogram; ranks 4 and 2 at \((5,1)\) & the same value test survives in a signed finite shadow \\
signed mixed layers & \(BI_3(0,1)\) has rank mass 300, below the size-\(B_3\) hash-ordered control 715 & signed incidence rows carry visible finite structure \\
dihedral \(m=8\) witness & same elementwise dihedral reflection-class histogram & centered standard ranks \(5\) and \(3\) detect relative mirror separation \\
dihedral \(m=9\) witness & same rational centered standard rank \(8\) & \(\SNFL\) separates \(1^8\) from \(1^6,2^2\) \\
length-colored \(G_2\) witness & short-axis and long-axis rows are ambient-conjugate after forgetting length & source profiles distinguish fixed short roots from fixed long roots \\
projective \(F_4\) witness & \(G_{\mathrm{proj}}\) and \(G_{\mathrm{mat}}\) act on the same 24 points but have centered ranks \(1\) and \(0\) & the admitted source layer matters before rank data is interpreted \\
\bottomrule
\end{tabular}
\end{table}

\begin{table}[!htbp]
\centering
\scriptsize
\setlength{\tabcolsep}{3pt}
\renewcommand{\arraystretch}{0.86}
\caption{What the worked witnesses show, II.}
\begin{tabular}{@{}L{0.24\textwidth}L{0.31\textwidth}L{0.35\textwidth}@{}}
\toprule
source & finite observation & allowed conclusion \\
\midrule
projective \(H_3\) witness & flat-size and Gram-angle channels act on the same 15 points with automorphism orders \(120\) and \(60\), while \((14,1)\) has rank \(0\) & the source channel can matter even when the ordinary standard block is blind \\
projective \(H_4\) witness & flat-size and Gram-square channels act on the same 60 projective root pairs with automorphism orders \(14400\) and \(7200\) & source-channel data can separate matroid/incidence symmetry from geometric Gram symmetry \\
projective \(D_5\) witness & pair channel has automorphism order \(122880\), while \(A_2\) incidence preservers have order \(1920\) and equal \(W(D_5)_{\mathrm{proj}}\) & source recovery may require triple/flat incidence, not just pair-level data \\
projective \(E_6\) witness & \(\mathrm{Par}_{012}\) and \(\mathrm{Par}_{0135}\) have size \(24\) and total rank \(8\), but \(U_{20}/U_{15}\) splits \((6,2)\) and \((5,3)\) & source-aware summand data can see parabolic row information hidden by total standard rank \\
projective \(E_7\) witness & \(\mathrm{Par}_{12345}\) and \(\mathrm{Par}_{12347}\) are both \(A_5\) rows of order \(720\), but have ranks \(5=(3,2)\) and \(6=(3,3)\) on \(U_{27}/U_{35}\) & same abstract parabolic type and order need not determine the embedded finite-shadow fingerprint \\
projective \(E_8\) witness & \(\mathrm{Par}_{1234567}\) and \(\mathrm{D8}_{1234567}\) are both \(A_7\) rows of order \(40320\), but have orbit profiles \([8,28,28,56]\) and \([1,28,28,28,35]\), with ranks \(3=(1,2)\) and \(4=(1,3)\) on \(U_{35}/U_{84}\) & same abstract subsystem type and order need not determine the source-aware finite-shadow fingerprint \\
M19/X05 exhibit & blind rank-locus mining repeatedly recovers the same order-72 group \(G_0\) when the full coset occurs & \(\Phi\) can recover a hidden group from centered-operator data \\
\texttt{lat565} X05 exception & local \(32+8+22\) centered-certificate stratification & \(\Phi\) separates a hidden-coset block from non-coset centered rank-drop certificates \\
\bottomrule
\end{tabular}
\end{table}

\subsection{A controlled X05 exhibit: M19 and \texttt{lat565}}

The X05 Sudoku/Latin row gives a separate controlled example of the same point.  For an order-nine square \(L\), take \(X=S_9\) acting by symbol relabeling and let
\[
  E_{\gamma,L}(i,j)=\gamma(L(i,j))-5
\]
be the centered operator, restricted to the standard subspace.  The fingerprint package records the rank-drop locus
\[
  \Gamma_{\mathrm{rank}}(L)=\{\gamma\in S_9:\operatorname{rank}E_{\gamma,L}<8\}
\]
together with the hidden stabilizer recovered by blind coset mining.  This is the same grammar slot \(\PhiFCIG\), now applied to a centered-operator sibling of the Type-A rows.

The published Sudoku operator paper isolates the original M19 rank-locus phenomenon; see \cite{BabanskyySudokuOperators2026}.  The centered Latin-square paper supplies the sibling determinant/rank-obstruction context for the Latin side of the centered-operator row\cite{BabanskyyLatin2026}.  The FCIG census source then asks whether M19 is isolated.  Its durable, source-locked answer is that the large order-72 hidden coset is not a one-off.  Whenever that full coset occurs in the corrected census, the hidden group is the same fixed subgroup \(G_0\leq S_9\), the affine symmetry group of the natural \(1,\ldots,9\) value grid.  In the corrected 1000+1000 census, this full \(G_0\)-coset appears in both Sudoku and Latin samples.  The visible magic-band refinement is observed only on the Sudoku side.  Complement relabeling is the baseline symmetry of the centered rank-drop condition.

The companion M19 census package now records the exceptional Latin square \texttt{lat565} as a local finite-certificate stratification
\[
  \Gamma_{\mathrm{rank}}(\texttt{lat565})=C_{32}\sqcup C_8\sqcup R_{22},\qquad 62=32+8+22.
\]
Here \(C_{32}\) is an order-32 paired-row hidden coset outside \(G_0\), \(C_8\) is a signed four-row parallelogram block, and \(R_{22}\) consists of eleven centered nullity-one complement-pair certificates.  Only \(C_{32}\) is group/coset-structured; the latter pieces are local centered certificates, not a second hidden group or a family theorem.

\begin{table}[!htbp]
\centering
\caption{M19/X05 census and \texttt{lat565} facts used here.}
\label{tab:m19-census}
\begin{tabular}{@{}lL{0.62\textwidth}@{}}
\toprule
item & recorded status \\
\midrule
full hidden coset & order 72; recurring signal, not a single M19 accident \\
recovered group & same fixed \(G_0\leq S_9\) whenever the full coset occurs \\
corrected census & 1000 Sudoku samples and 1000 Latin samples inspected in the source package \\
magic-band refinement & observed on the Sudoku side only \\
\texttt{lat565} exception & local \(32+8+22\) stratification; only the 32-point block is a hidden coset; the 8-point and 22-point pieces are centered linear certificates only \\
\bottomrule
\end{tabular}
\end{table}

This exhibit is useful because \(\PhiFCIG\) recovers a recurring finite group from the data without being handed the group in advance, and because \texttt{lat565} shows the same channel separating a hidden-coset block from non-coset centered certificates inside one object.  It is not used here as a theorem pillar.  The frequencies are empirical, the mechanism selecting the arrays is open, and the broader out-of-\(G_0\) population remains outside this paper.

\subsection{Source quotients require coupling data}
\label{subsec:source-coupling}

The quotient mechanism is classical.  Higman fixes the arbitrary commutative coefficient ring in Section~6 and gives the unnormalised block-row-sum quotient in Section~8, Eq.~(8.1)\cite[\S6, p.~11; \S8, Eq.~(8.1), p.~20]{Higman1970CoherentConfigurations}; Godsil records the equitable-partition identity \(AH=HB\), the invariant cell-constant subspace, and spectral divisibility\cite[\S5.1, Eq.~(5.1.1) and Lem.~5.1.2]{GodsilAssociationSchemes2010}; and Haemers supplies the classical quotient/interlacing context\cite[\S2]{Haemers1995Interlacing}.  This paper uses the mechanism only as an interface guardrail showing which source and coupling data must be retained; it makes no novelty claim for quotient existence, spectral inheritance, or interlacing.

Several rows above compare a row with a quotient row.  A block partition \(\Omega=\bigsqcup_B\Omega_B\) compatible with a centered operator \(N\) (constant block-column sums) induces a quotient operator \(N_\pi\) on the blocks, and, after choosing one base point \(b(B)\in\Omega_B\) per block, the anchored augmentation lattice splits into a kernel lattice and a lifted quotient lattice, over which \(N\) becomes block triangular,
\[
  N\;\sim\;\begin{pmatrix} N_K & \beta \\ 0 & N_\pi \end{pmatrix},
\]
where \(N_K\) is the restriction to the kernel lattice and \(N_\pi\) here denotes the quotient operator on its own anchored lattice.  The grammar question is when the two endpoint fingerprints determine the total one, and the honest answer is a guardrail: only together with the coupling \(\beta\).

The underlying descent bookkeeping can be stated without normalized averaging.
Let \(q:\mathbb Z[\Omega]\to\mathbb Z[\pi]\) send every point to its
block and write \(A_X\) for the augmentation lattice of a finite set \(X\).
Then
\[
  0\longrightarrow K_\pi:=\bigoplus_{B\in\pi}A_{\Omega_B}
  \longrightarrow A_\Omega\xrightarrow{q}A_\pi\longrightarrow0
\]
is split exact after base points are chosen, and an operator descends exactly
when
\[
  qN=N_\pi q
  \quad\Longleftrightarrow\quad
  N(\ker q)\subseteq\ker q.
\]
These are the classical quotient/equitable conditions in augmentation
language.  For a section \(s:A_\pi\to A_\Omega\), put
\[
  \beta_s=N_As-sN_{\pi A}.
\]
If \(s_h=s+i h\), where \(i:K_\pi\hookrightarrow A_\Omega\), direct block
multiplication gives
\[
  \beta_{s_h}=\beta_s+N_Kh-hN_{\pi A}.
\]
Consequently, for \(v\in\ker N_{\pi A}\), the connecting class
\[
  \delta(v)=[\beta_s(v)]\in\operatorname{coker}N_K
\]
is independent of the chosen section.  This class measures one obstruction,
but it does not classify integral equivariant complements.

The modular-centered row \(I_7(1,1)\) supplies the useful falsifier.  In its
integral split augmentation basis,
\[
  N_K=0_3,
  \qquad
  N_A=\begin{pmatrix}0&X\\0&2X\end{pmatrix},
  \qquad
  X=\begin{pmatrix}0&-4&-4\\-4&0&-4\\-4&-4&0\end{pmatrix}.
\]
Here \(\delta_{\mathbb Z}=0\), but the section equation
\(\beta_s+N_Kh-hN_{\pi A}=0\) has only \(h=\frac12I_3\) over
\(\mathbb Q\) and no integral solution.  Thus a zero connecting map does not
force an integral equivariant complement.  Independently, elementary
left--right unimodular operations reduce \(N_A\) to \(X\oplus0_3\); the
determinantal divisors \(4,16,128\) recover the nonzero anchored Smith factors
\([4,4,8]\) from the \(3\times3\) block, without a full ambient Smith
reduction.

\begin{proposition}[source-side splitting test]\label{prop:source-coupling}
For each block \(B\) write
\[
  d(B)=N e_{b(B)}-\sum_{B'}(N_\pi)_{B'B}\,e_{b(B')},
\]
a vector supported in the kernel lattice and computable from the base-point columns of \(N\) and the quotient matrix alone, with no Smith computation.  If all the \(d(B)\) coincide, then \(\beta=0\) for the base-point section, \(N\) is conjugate over \(\mathbb Z\) to \(N_K\oplus N_\pi\) on the anchored lattice by a unitriangular change of section, and \(\SNFL(N)\) is the Smith form of the direct sum: the multiset of elementary divisors is the union of the endpoint multisets, with the invariant factors re-chained by gcd/lcm.
\end{proposition}

\begin{proof}
On the lifted basis vector \(h_C=e_{b(C)}-e_{b(C_0)}\) a direct computation gives \(\beta(h_C)=d(C)-d(C_0)\), so equal defects force \(\beta=0\); the splitting statement is then block-diagonal bookkeeping, and the Smith form of a direct sum is classical.
\end{proof}

The hypothesis is not removable as a general source-side test, and endpoint data alone never suffices.  The three-point row
\(N=\begin{psmallmatrix}3&1&0\\-1&1&0\\-2&-2&0\end{psmallmatrix}\)
with blocks \(\{1,2\}\) and \(\{3\}\) has both anchored endpoint operators equal to \((2)\), yet its anchored total operator is \(\begin{psmallmatrix}2&1\\0&2\end{psmallmatrix}\) with \(\SNFL=[1,4]\): the total cokernel is \(\mathbb Z/4\), while the uncoupled direct sum gives \([2,2]\), that is \(\mathbb Z/2\oplus\mathbb Z/2\).  The missing datum is precisely the extension class of the coupling.

For completeness, the residue description used here has an elementary proof.  After endpoint Smith transforms write the block matrix as
\[
  \begin{pmatrix}A&C\\0&B\end{pmatrix},
  \qquad A=\operatorname{diag}(a_i),\quad B=\operatorname{diag}(b_j).
\]
Endpoint-preserving block-unitriangular row and column operations replace \(C\) by \(C+YB+AZ\).  Entrywise, this replaces \(c_{ij}\) by \(c_{ij}+y_{ij}b_j+a_i z_{ij}\), so the invariant is
\[
  [c_{ij}]\in \mathbb Z/(a_i,b_j)
  \;\cong\;
  \mathbb Z/\gcd(a_i,b_j)\mathbb Z.
\]
Conversely, B\'ezout's identity shows entrywise that two couplings with the same residues differ by such a pair \(Y,Z\).  Thus the residue matrix is exactly the coupling class relative to the chosen endpoint Smith bases, not a new structure theorem.

Deciding splitting in general is Sylvester-equation solvability \cite{Roth1952,Gustafson1979}; when the endpoint determinants are coprime the coupling never matters \cite[Thm.~2]{Newman1974}; and the possible middle types at a shared prime form the classical Green--Klein/Littlewood--Richardson set \cite{Klein1968}; see also \cite[Thms.~4.1--4.2]{Schmidmeier2007LR}.  This subsection is therefore an interface guardrail, not a new descent or transfer theory: no shared-prime law and no novelty are claimed.

The same audits also fix the limit.  The fingerprint is not a complete invariant, rational rank is not the whole fingerprint, and rank data is not used as a decision oracle for unrelated structural questions.  This section records the value and the guardrails together.
